\section{Accumulator SPIN}
\label{sec:accumulator-spin}

Accumulator SPIN is the transparent instantiation of the framework. It uses a
random block-diagonal outer, a uniform interleaver, and the binary
accumulator.

\subsection{Relation to Block-Accumulate codes}

Block-Accumulate (BA) codes concatenate a random block-diagonal outer with
one or more stages, each consisting of a uniform interleaver and an accumulator
\cite{KolesnikovEtAl2026BA}. Accumulator SPIN retains the first such stage and
therefore may be viewed as a truncated BA code. We retain the SPIN name to
place this construction in the common outer--interleaver--inner interface, but
the block-accumulate architecture and the random-block enumerator below are
due to Kolesnikov et al.

Accumulator SPIN follows the BA ensemble at this level: setup samples the
local maps independently. The theorem below conditions each local map to be
injective. This condition makes every realized outer encoder injective.

Chosen-block BAA studies the related fixed-constituent choice
\cite{PecenyRindal2026ChosenBlockBAA}. It fixes a short constituent code and
uses the same code in every outer position. Its concrete binary constituents
include Golay, Reed--Muller, and a BCH subcode. That reuse option is not the
ensemble analyzed here.

\subsection{Random-block outer lineage}

The following local input--output enumerator is the random-block calculation
used for BA codes. We state it in the present notation because it is also the
local input to the Accumulator SPIN analysis.

\begin{lemma}[Uniform random-block enumerator]
\label{lem:uniform-random-block-enumerator}
Suppose that $G_B$ is uniform over $\F_2^{K_B\times B}$. Let
$E^{\mathrm{blk}}_{a,b}(G_B)$ be the number of inputs of weight $a$ whose
images have weight $b$. Then
\begin{equation}
  \E_{G_B}[E^{\mathrm{blk}}_{a,b}(G_B)]
  =
  \begin{cases}
    1, & a=0 \text{ and } b=0,\\
    0, & a=0 \text{ and } b>0,\\
    \displaystyle
    \binom{K_B}{a}\binom{B}{b}2^{-B}, & a>0.
  \end{cases}
  \label{eq:uniform-random-block-enumerator}
\end{equation}
\end{lemma}

\begin{proof}
The zero input maps to the zero output. Fix a nonzero input $u$. Each column
of $G_B$ is uniform and independent, so the coordinates of $uG_B$ are
independent uniform bits. Thus $uG_B$ is uniform over $\F_2^B$. There are
$\binom{K_B}{a}$ inputs of weight $a$ and $\binom{B}{b}$ outputs of weight
$b$. Summing the probability $2^{-B}$ over these pairs gives
\eqref{eq:uniform-random-block-enumerator}.
\end{proof}

The BA independent-block ensemble has a closed global input--output
enumerator. Let $E^{\mathrm{BA}}_{a,b}$ count total input weight
$a$ and total outer-output weight $b$ across $M$ blocks. Specializing the BA
formula~\cite{KolesnikovEtAl2026BA} to local dimensions
$K_B\times B$ gives
\begin{equation}
  \E[E^{\mathrm{BA}}_{a,b}]
  =
  \sum_{q=0}^{M}
  2^{-Bq}\binom{M}{q}\binom{Bq}{b}
  \sum_{j=0}^{q}(-1)^j\binom{q}{j}
  \binom{K_B(q-j)}{a}.
  \label{eq:ba-independent-block-enumerator}
\end{equation}
Here $q$ is the number of nonzero input blocks. The inner alternating sum
counts weight-$a$ inputs for which exactly those $q$ blocks are nonzero. For a
fixed such input, the $q$ independently sampled local matrices produce a
uniform word on the corresponding $Bq$ output coordinates. This yields the
factor $2^{-Bq}\binom{Bq}{b}$.

The theorem below uses the same independence pattern but samples uniform
injections rather than unconditioned matrices. Its exact generating function
appears in Section~\ref{sec:accumulator-spin-theorem}.

\subsection{Accumulator inner}

Define $\Acc_N:\F_2^N\to\F_2^N$ by
\begin{equation}
  (\Acc_N(u))_t
  :=
  \sum_{i=1}^t u_i
  \pmod 2,
  \qquad t\in[N].
  \label{eq:accumulator-definition}
\end{equation}
Equivalently, if $y:=\Acc_N(u)$ and $y_0:=0$, then
$y_t=y_{t-1}\mathbin\oplus u_t$.

For $w,h\in\{0,\ldots,N\}$, define
\[
  A^{\Acc}_{w,h}
  :=
  |\{u\in\F_2^N:\wt(u)=w,\ \wt(\Acc_N(u))=h\}|.
\]

\begin{lemma}[Exact accumulator enumerator]
\label{lem:accumulator-enumerator}
Fix $w\in\{1,\ldots,N-1\}$ and set $k:=\lceil w/2\rceil$. Then
\begin{equation}
  A^{\Acc}_{w,h}
  =
  \binom{h-1}{k-1}\binom{N-h}{w-k}.
  \label{eq:accumulator-enumerator}
\end{equation}
\end{lemma}

\begin{proof}
Let the ones of $u$ occur at
$1\le p_1<\cdots<p_w\le N$. Set $p_0:=0$, $p_{w+1}:=N+1$, and
$\Delta_i:=p_i-p_{i-1}$. The positive spacings satisfy
$\sum_{i=1}^{w+1}\Delta_i=N+1$.

The accumulator output equals one between its first and second flips, between
its third and fourth flips, and so on. Hence its weight is the sum of the
even-index spacings. There are $k$ such spacings. The number of positive
compositions of $h$ into those $k$ parts is
$\binom{h-1}{k-1}$. The remaining $w+1-k$ spacings sum to $N+1-h$,
which gives $\binom{N-h}{w-k}$ choices. Multiplication proves
\eqref{eq:accumulator-enumerator}.
\end{proof}

If $U_w$ is uniform over the weight-$w$ slice, then
\begin{equation}
  \Prob[\wt(\Acc_N(U_w))\le d]
  =
  \frac{1}{\binom Nw}
  \sum_{h=0}^d A^{\Acc}_{w,h}.
  \label{eq:accumulator-exact-tail}
\end{equation}

\begin{lemma}[A simple accumulator tail]
\label{lem:accumulator-tail}
Fix $\delta\in(0,1/4)$ and
$w\in\{1,\ldots,\lfloor N/2\rfloor\}$. Then
\[
  \Prob[\wt(\Acc_N(U_w))\le\lfloor\delta N\rfloor]
  \le
  (4e\delta)^{\lceil w/2\rceil}.
\]
\end{lemma}

\begin{proof}
Set $k:=\lceil w/2\rceil$. Lemma~\ref{lem:accumulator-enumerator} and the
hockey-stick identity give
\[
  \sum_{h\le\lfloor\delta N\rfloor}A^{\Acc}_{w,h}
  \le
  \binom{N}{w-k}\binom{\lfloor\delta N\rfloor}{k}.
\]
After division by $\binom Nw$, bound the binomial term by
$(e\delta N/k)^k$. Since $w\le N/2$, the remaining binomial ratio is at
most $(2w/N)^k$. Finally, $w\le2k$, which gives the claim.
\end{proof}

The next contraction is uniform in the input weight. It is the form used by
the asymptotic theorem.

\begin{lemma}[Sharp accumulator contraction]
\label{lem:sharp-accumulator-contraction}
Fix $\delta\in(0,1/2)$ and define
\[
  \rho(\delta):=2\sqrt{\delta(1-\delta)}.
\]
For every $w\in\{1,\ldots,N\}$,
\begin{equation}
  \Prob[\wt(\Acc_N(U_w))\le\lfloor\delta N\rfloor]
  \le \rho(\delta)^w.
  \label{eq:sharp-accumulator-contraction}
\end{equation}
\end{lemma}

\begin{proof}
Set $D:=\lfloor\delta N\rfloor$ and $k:=\lceil w/2\rceil$.
For $w<N$, Lemma~\ref{lem:accumulator-enumerator} identifies the probability
in~\eqref{eq:sharp-accumulator-contraction} with $\Prob[X\ge k]$, where $X$
is hypergeometric with sample size $w$ and $D$ marked positions among $N$.
For $t>0$, the moment bound for sampling without replacement gives
\[
  \E[e^{tX}]
  \le (1-\delta+\delta e^t)^w.
\]
Markov's inequality and $k\ge w/2$ give
\[
  \Prob[X\ge k]
  \le
  \inf_{t>0}
  \left(e^{-t/2}(1-\delta+\delta e^t)\right)^w.
\]
The minimizer satisfies $e^t=(1-\delta)/\delta$, which yields the claim.
When $w=N$, the accumulator output has weight $\lceil N/2\rceil>D$, so the
left side is zero.
\end{proof}

\subsection{Asymptotic theorem}
\label{sec:accumulator-spin-theorem}

Fix a rational
rate $R\in(0,1)$ and let
\[
  \mathcal B_R:=\{B\ge1:RB\in\mathbb Z\}.
\]
For $B\in\mathcal B_R$ and each $i\in[M]$, setup samples an independent
uniform linear injection
\[
  C_i:\F_2^{RB}\longrightarrow\F_2^B.
\]
It then samples $\Pi\gets S_N$, where $N:=MB$. For
$x=(x^{(1)},\ldots,x^{(M)})$, define
\begin{equation}
  E_N^{\mathrm{acc}}(x)
  :=
  \Acc_N\!\left(
    \Pi(C_1(x^{(1)})\Vert\cdots\Vert C_M(x^{(M)}))
  \right).
  \label{eq:independent-accumulator-spin}
\end{equation}
The setup remains fixed for all messages. Every realized encoder is
injective.

For one local injection, define
\begin{equation}
  G_{B,R}(z)
  :=
  (2^{RB}-1)\frac{(1+z)^B-1}{2^B-1}.
  \label{eq:injection-local-mass}
\end{equation}
Every nonzero input has a uniform nonzero image, so independence across block
positions gives
\begin{equation}
  \E\!\left[
    1+\sum_{w=1}^{N}A_{N,w}^{\mathrm{out}}z^w
  \right]
  =(1+G_{B,R}(z))^M.
  \label{eq:injection-global-mass}
\end{equation}

Define
\begin{equation}
  \lambda_A(R,\delta)
  :=
  1-R-\log_2(1+\rho(\delta)).
  \label{eq:accumulator-sparse-exponent}
\end{equation}
Fix $c>0$. For each sufficiently large $M$, let $B_M$ be the least member of
$\mathcal B_R$ satisfying
\[
  B_M\ge c\log_2(MB_M),
\]
and set $N_M:=MB_M$.

\begin{theorem}[Accumulator SPIN]
\label{thm:independent-accumulator-spin}
Suppose
\[
  \lambda_A(R,\delta)>0,
  \qquad
  c\lambda_A(R,\delta)\ge1.
\]
Then
\begin{equation}
  \Prob[d_{\min}(E_{N_M}^{\mathrm{acc}})\le\delta N_M]
  =o(1)
  \qquad(M\to\infty).
  \label{eq:independent-accumulator-distance}
\end{equation}
The family has rate $R$, block size $\Theta(\log N_M)$, and
$\Theta(N_M\log N_M)$ direct encoding cost.
\end{theorem}

\begin{proof}
Let $D:=\lfloor\delta N_M\rfloor$ and let $Z_D$ count nonzero messages whose
encoded words have weight at most $D$. Corollary~\ref{cor:weight-first-moment},
Lemma~\ref{lem:sharp-accumulator-contraction}, and
\eqref{eq:injection-global-mass} give
\[
  \E[Z_D]
  \le
  (1+G_{B_M,R}(\rho(\delta)))^M-1.
\]
For $z>0$, equation~\eqref{eq:injection-local-mass} gives
\[
  G_{B,R}(z)
  \le
  2\left(2^{R-1}(1+z)\right)^B.
\]
Set $\lambda:=\lambda_A(R,\delta)$. The block schedule implies
\[
  M G_{B_M,R}(\rho(\delta))
  \le
  \frac{2}{B_M}N_M^{1-c\lambda}
  =o(1).
\]
Thus $\E[Z_D]=o(1)$, and Markov's inequality proves
\eqref{eq:independent-accumulator-distance}.
\end{proof}

At rate $R=1/2$, the choice $c=3$ and
$\delta=37/10000$ satisfies the theorem with a strict interval-certified
margin. The corresponding relative-distance guarantee is small but positive.
For this proof, the largest possible rate-half limit as $c\to\infty$ is
$0.0449101394\ldots$.
